\documentclass[letterpaper,11pt]{article}

% \usepackage[style=ieee,url=false,doi=false,maxbibnames=99,sorting=ydnt,dashed=false,backend=biber]{biblatex}

\usepackage[style=ieee,url=true,doi=true,maxbibnames=99,sorting=none,dashed=false,backend=biber]{biblatex}

\usepackage{geometry}  
\geometry{a4paper,left=25mm,right=20mm,top=40mm,bottom=25mm}
\usepackage{pageslts,refcount}
\usepackage{lastpage}
\usepackage{fancyhdr}
\pagestyle{fancy}
\bibliography{references_cv}
\usepackage{simplecv}

\cfoot{\thepage\ of \pageref{LastPage}}               
\rfoot{\ifnum\getpagerefnumber{VeryLastPage}=\value{page} \textit{Last Updated: \today} \fi}

\begin{document}

% Heading
\headinginline{Saurabh Kumar}{
    Google Scholar: \website{tinyurl.com/4pyub4pd} \\
    GitHub: \github{saurabhk0317} \\
    Email: \email{saurabhk0317@gmail.com} \\
    LinkedIn: \linkedin{saurabhk0317}
}

\section{Education}
\textbf{Indian Institute of Science (IISc)} \hfill Bengaluru, India\\
M.Tech. (Research) in Electrical Engineering \hfill 2026\textendash Present\\
Advisor: Prof. Prasanta Kumar Ghosh (SPIRE Lab) \\
GATE 2026 (Data Science and AI): AIR 362/69{,}242
Publications: \cite{kumar2026wsasr}
\begin{itemize}
  \item Proposed a CTC training criterion with token-level wildcard arcs that tolerates localized transcription ambiguity (pronunciation, spelling, lexical variants), achieving a \textbf{9.45\% average relative WER reduction over CTC} across multiple languages and corpora~\cite{kumar2026wsasr}.
\end{itemize}

\medskip
\textbf{National Institute of Technology, Patna} \hfill Patna, India\\
B.Tech. in Electronics and Communication Engineering \hfill 2016\textendash 2020 \\
GPA: 7.42/10.00

\section{Experience}
\textbf{SPIRE Lab, Indian Institute of Science (IISc)} \hfill Bengaluru, India\\
Research Associate \hfill July 2021\textendash July 2026 \\
Advisor: Prof. Prasanta Kumar Ghosh \\
Projects:
\begin{itemize}
  \item \href{https://respin.iisc.ac.in/}{\textbf{RESPIN}}: \textit{Speech Recognition in Agriculture and Finance for the Poor in India} \\
  Publications: \cite{kumar2025respins, ASRU25_MADASR2, Saurabh_IS25, ICASSP25_DID, IS24_ASR, ASRU23_RESPIN}
  \begin{itemize}
    \item First author of the RESPIN-S1.0 database paper published at \textbf{NeurIPS 2025 (Datasets and Benchmarks Track)}; led large-scale data curation, reproducible benchmarking, and baseline ASR development across nine Indian languages~\cite{kumar2025respins}.
    
    \item Developed a multimodal fusion model for joint ASR and dialect identification (DID) (\textbf{4.65\% CER and 81.63\% DID accuracy}), integrating speech and text embeddings for Indian dialects~\cite{Saurabh_IS25}.
    
    \item Led the first large-scale study on DID in Indian languages (\textbf{79.81\% accuracy across 33 dialects})~\cite{ICASSP25_DID}.
    
    \item Enhanced low-resource ASR via adapter-based domain adaptation, reducing \textbf{WER (10\%-30\%)} and \textbf{CER (3\%-18\%)} across 9 languages~\cite{IS24_ASR}.
    
    \item Proposed a Gated Multi-Encoder ASR model with frame-level gating, reducing \textbf{WER by 7.5\% (Bhojpuri) and 9\% (Bengali)}~\cite{ASRU23_RESPIN}.
  \end{itemize}

  \item \href{https://syspin.iisc.ac.in/}{\textbf{SYSPIN}}: \textit{Text-to-Speech Synthesizer in Nine Indian Languages} \\
  Publications: \cite{LIMMITS24_OJSP2025}
  \begin{itemize}
    \item Developed a data filtering pipeline using Kaldi-based ASR models, improving corpus quality.
  \end{itemize}
\end{itemize}

\textbf{ARTPARK @ IISc} \hfill Bengaluru, India\\
Machine Learning/Signal Processing Consulting Engineer \hfill Sep 2022\textendash July 2026 \\
Project: 
\begin{itemize}
\item \href{https://vaani.iisc.ac.in/}{\textbf{VAANI}}: \textit{Capturing the language landscape for an inclusive digital India}
\begin{itemize}
  \item Developed automated signal processing and machine learning pipelines for acoustic quality assessment, speech activity detection, and metadata validation.
  \item Established quantitative acceptance criteria and automated corpus curation protocols underpinning the open-source release of \textbf{31{,}000+ hours} of speech from \textbf{156{,}000+ speakers} across \textbf{109 languages} and \textbf{31 states and union territories}.
\end{itemize}
\end{itemize}

\textbf{Department of ECE, NIT Patna} \hfill Patna, India\\
Undergraduate Student Researcher \hfill July 2019\textendash June 2021 \\
Advisor: Dr. S. Shahnawazuddin \\
Publications: ~\cite{Syed2021_APAC, Syed2021_DSP}
\begin{itemize}
  \item Investigated formant-modification and bispectrum-based front-end features to improve children’s ASR under zero-resource conditions~\cite{Syed2021_APAC}.
  \item Developed a CGAN-based voice conversion framework, transforming adult speech into child-like speech for ASR adaptation~\cite{Syed2021_DSP}.
\end{itemize}

\textbf{EICT Academy, IIT Guwahati} \hfill Guwahati, India\\
Summer Intern \hfill May 2019\textendash June 2019 \\
\begin{itemize}
  \item Designed Verilog code for a 5-stage discrete wavelet transform cascade filter bank for FPGA-based simulation.
\end{itemize}

\section{Mentorship \& Professional Services}

\textbf{Reviewer:}
\begin{itemize}
    \item \textbf{Conferences and Workshops:} ICASSP (since 2025), IJCNN (since 2025), IEEE ASRU 2025, Interspeech (since 2026), IEEE SLT 2026.
    \item \textbf{Journals:} Computer Speech \& Language (Elsevier), Speech Communication (Elsevier).
\end{itemize}

\textbf{Organizer:}  
\begin{itemize}
    \item \textbf{ASRU 2025 (MADASR 2.0):} Lead organizer of the Multi-Lingual Multi-Dialect ASR Challenge in 8 Indian Languages; designed challenge tracks, evaluation protocols, public leaderboards, developed ESPnet-based baseline systems, and reviewed challenge track submissions.
    \item \textbf{ICASSP 2024 (LIMMITS’24):} Co-organized session on Indic TTS, developed baselines, and reviewed challenge track submissions.
    \item \textbf{ASRU 2023 (MADASR):} Led data curation, ASR baselines for Bengali, Bhojpuri.
    \item \textbf{Interspeech 2022 (Gram Vaani ASR):} Organized Hindi ASR challenge, developed Kaldi baselines.
\end{itemize}

\textbf{Student Volunteer:} ICASSP 2024, ICASSP 2025.  
  

\section{Awards}
\textbf{SLT-CODE Hackathon Winner 2022} – Won \textbf{Best Hackathon Project} Award for developing dialectal ASR systems for Bengali and Bhojpuri, presented at \textbf{SLT 2022}.


\printbibliography[title=\textbf{Selected Publications}]

\section{Languages and Skills}
\begin{minipage}{0.48\textwidth}
    \textbf{Languages}
    \begin{itemize}
        \item \textbf{Programming:} Python, Bash, PHP
        \item \textbf{Natural:} Bhojpuri, English, Hindi, Maithili
        % \item \textbf{TOEFL:} 94 (Reading: 23, Listening: 26, Speaking: 23, Writing: 22)
    \end{itemize}
\end{minipage}
\hfill
\begin{minipage}{0.48\textwidth}
    \textbf{Skills}
    \begin{itemize}
        \item \textbf{ML/DL Toolkits:} PyTorch, Scikit-learn
        \item \textbf{ASR Frameworks:} ESPnet, Kaldi, Fairseq
        \item \textbf{Others:} MATLAB, Audacity, Praat, Git
    \end{itemize}
\end{minipage}

\end{document}
